\documentclass[11pt]{article}
% Running overview for colleagues (informal). For two columns: add "twocolumn" to the class options.
\usepackage[margin=1in]{geometry}
\usepackage{amsmath,amssymb}
\usepackage{booktabs}
\usepackage{float}
\usepackage{microtype}
\usepackage{xcolor}
\usepackage[authoryear,round]{natbib}
\usepackage[colorlinks=true,linkcolor=blue!50!black,citecolor=blue!50!black,urlcolor=blue!50!black]{hyperref}
\setlength{\parskip}{0.5em}
\setlength{\parindent}{0pt}

\title{\bfseries When does a quantum feature map do something classical methods can't copy?\\[0.3em]
\large A running overview of the teacher--student experiments}
\author{Christopher Clark}
\date{Draft, \today}

\begin{document}
\maketitle

% =====================================================================
\section{What we're doing and why}

\textbf{What we're doing and why.} Quantum feature maps take each sample, use it to steer a small quantum
system, and read out a set of measurements as new features. The hope is that some of those features are hard
to get any other way. Several papers report that these maps beat classical models on biomedical data
\citep[e.g.][]{Simen2025}. But when we tried to reproduce those results, ordinary classical models matched them,
and our own run on 55 datasets showed no general advantage. The trouble is that a loss on real data is
ambiguous: maybe the quantum map can't produce anything special, or maybe those datasets just didn't need it.
So we flipped the question around. We let the quantum map \emph{make} the labels itself, then checked who can
learn them: the full quantum simulation, cheaper non-quantum versions of the same map (one that averages the
qubits' influence on each other, and one with no interaction at all), and standard classical models. Doing
this across the map's settings (coupling strength $J$, anneal time $\tau$, which measurements we read out, and
the number of qubits) tells us where, if anywhere, the quantum part is doing something nothing else can copy.
That's also where it would be worth looking for real data, and spending hardware time.

\textbf{Why it might work.} There's good reason to think quantum systems \emph{can} produce patterns that
classical methods struggle with. Researchers have built carefully designed examples where a quantum learner
provably beats any classical one \citep{Liu2021}. The physics community has also argued for years that
systems of interacting qubits, much like ours, make rich, useful features for machine learning
\citep{Fujii2017}. And our own first experiment found one case where neither classical models nor the cheap
non-quantum versions of the map could keep up.

\textbf{Why it might not.} The field's recent track record leans the other way. A lot of quantum machine
learning models turn out to be built from fairly simple ingredients, essentially sums of waves
\citep{Schuld2021}, and classical tricks can often mimic them closely
\citep{Schreiber2023,Landman2022,Sweke2025}. With enough training data, classical learners often catch up
even on labels a quantum system produced \citep{Huang2021}. Careful head-to-head benchmarks usually find
classical models do just as well \citep{Bowles2024}, and models that are easy to train tend to be the ones
that are also easy to simulate classically \citep{Cerezo2023}. Two things specific to our setup also cut
against us. First, we couple every qubit to every other one, which is exactly the case where the cheap
``averaged'' version gets \emph{more} accurate as the system grows, so any advantage could shrink at larger
sizes. Second, at the sizes we can test (up to 8 qubits) the quantum simulation itself runs on an ordinary
computer, so even a clear win here would show that the map's \emph{structure} fits certain data well, not that
a quantum computer is needed.

% =====================================================================
\section{How the data are made}

Each dataset is a table of random numbers plus a 0/1 label for every row. The numbers themselves carry no
meaning; all the structure is in how the labels are made. We run the quantum map (or one of its classical
versions) on every row, combine its measurements into a score, and turn the score into a label.

\paragraph{Step 1: inputs.} We draw $n=600$ samples, each with $p$ features, where $p$ is also the number of
qubits (one qubit per feature). For $p=4$ the inputs are a $600\times4$ matrix of random numbers, one row per
sample and one column per qubit:
\begin{equation}
x^{(m)}\sim\mathcal N(0,I_p),\qquad m=1,\dots,600 .
\end{equation}
Each column is standardized and clipped at $\pm3$:
$z_k=\mathrm{clip}\big((x_k-\bar x_k)/\mathrm{sd}(x_k),\,-3,\,3\big)$.
In the main run, every dataset with the same $p$ uses the \emph{same} 600 inputs, so comparisons across
settings are paired.

\paragraph{Step 2: the map.} Each sample sets the fields on the $p$ qubits. Starting from all qubits pointing
along $+x$, $|\psi(0)\rangle=|+\rangle^{\otimes p}$, we anneal over a time $\tau$:
\begin{equation}
H(s)=(1-s)\Big(-\sum_{i=1}^{p}X_i\Big)+s\Big(-\sum_{i=1}^{p}z_iZ_i-J\sum_{i<j}Z_iZ_j\Big),\qquad s=t/\tau\in[0,1].
\end{equation}
The first term is the starting ``driver'' and fades out. The second is the data Hamiltonian and fades in:
feature $z_i$ pushes qubit $i$ up or down, and $J$ makes every pair of qubits prefer to point the same way.
At the end ($s=1$) we record, for every qubit and every pair,
\begin{equation}
\langle X_i\rangle,\ \langle Y_i\rangle,\ \langle Z_i\rangle,\qquad
\langle X_iX_j\rangle,\ \langle Y_iY_j\rangle,\ \langle Z_iZ_j\rangle .
\end{equation}

\paragraph{Step 3: three ways to compute the same anneal (the ``teachers'').}
\begin{itemize}\setlength{\itemsep}{2pt}
\item \textbf{Exact (Q):} the full quantum state ($2^p$ numbers), evolved by the Schr\"odinger
equation.\footnote{Numerically: $N=\max(400,\lceil160\tau\rceil)$ time steps; each step applies half of the
data term (as phases), the driver (as rotations of every qubit), then the other half of the data term
(Strang splitting).}
\item \textbf{Mean-field (MF):} each qubit evolves on its own, feeling its own feature plus the
\emph{average} pull of the others,
\begin{equation}
H_i(t)=-(1-s)X_i-s\,h_i(t)\,Z_i,\qquad h_i(t)=z_i+J\sum_{j\ne i}\langle Z_j\rangle_t ,
\end{equation}
with the averages updated at every step and pair terms taken as products,
$\langle X_iX_j\rangle=\langle X_i\rangle\langle X_j\rangle$. No entanglement.
\item \textbf{No coupling ($J=0$):} as mean-field with $h_i=z_i$. Each qubit anneals alone, so every feature
depends on one input only.
\end{itemize}

\paragraph{Step 4: the settings.} For each $p$ we set $J=c/(p-1)$ with $c\in\{0.5,1,2,4,8\}$
(Table~\ref{tab:J}).
\begin{table}[H]\centering\small
\caption{Coupling values $J$ used at each $p$.}\label{tab:J}
\begin{tabular}{@{}cccc@{}}\toprule
$c$ & $p=4$ & $p=6$ & $p=8$\\\midrule
0.5 & 0.167 & 0.100 & 0.071\\
1 & 0.333 & 0.200 & 0.143\\
2 & 0.667 & 0.400 & 0.286\\
4 & 1.333 & 0.800 & 0.571\\
8 & 2.667 & 1.600 & 1.143\\\bottomrule
\end{tabular}
\end{table}

\begin{center}\small
\begin{tabular}{@{}ll@{}}\toprule
setting & values\\\midrule
qubits $p$ & 4, 6, 8\\
coupling $J$ & Table~\ref{tab:J}\\
anneal time $\tau$ & 0.5, 1, 1.5, 2, 3, 4, 5, 7, 10, 15, 20, 40 \quad (earlier run: 1, 5)\\
readout & \textbf{full}: all $\langle X_i\rangle,\langle Y_i\rangle,\langle Z_i\rangle$ ($3p$ numbers);\\
        & \textbf{$Z$-only}: $\langle Z_i\rangle$ ($p$ numbers; what an annealer measures)\\\bottomrule
\end{tabular}
\end{center}
(The earlier run also tried a pair-term readout and the single sum $\sum_i\langle X_i\rangle$.)

\paragraph{Step 5: from measurements to a label.} Collect the readout into a matrix $F$, one row per sample.
Standardize each column, draw random weights, and form a score:
\begin{equation}
\tilde F_k=\frac{F_k-\bar F_k}{\mathrm{sd}(F_k)},\qquad w\sim\mathcal N(0,I),\qquad
\mathrm{score}=\mathrm{standardize}\big(\tilde F w\big).
\end{equation}
Then flip a weighted coin:
\begin{equation}
y\sim\mathrm{Bernoulli}\big(\sigma(\mathrm{score}/T)\big),\qquad \sigma(u)=\frac{1}{1+e^{-u}},\quad T=0.25 .
\end{equation}
The random weights make the label a \emph{generic} function of the map's output rather than one we picked.
The temperature $T$ makes the labels noisy, as in real data: even a model that knew the true score would only
reach AUC $\approx0.95$. Weights and coin flips are seeded by the dataset's name, so every dataset can be
regenerated exactly.

\paragraph{Step 6: how many.} For each ($p$, $\tau$, readout) we make 12 datasets: 2 with no coupling (two
seeds), plus one mean-field and one exact for each of the 5 couplings. That is
$3\times12\times2\times12=864$ datasets in the main run (216 in the earlier one). The learners only ever see
$(x,y)$, i.e.\ the $600\times p$ input matrix and the 600 labels.

% =====================================================================
\section{Who tries to learn the labels (the students)}

Every student sees the same thing: the $600\times p$ inputs and the 600 labels. Each one turns the inputs into
features in its own way and fits a \textbf{logistic regression} on them (the nonlinear learners use their own
model instead). The students differ only in \emph{which} features they are allowed to use.

\begin{table}[H]\centering\small
\caption{The students and their menus. The 9 measurement groups are $X$, $Y$, $Z$, all three, $XX$, $YY$,
$ZZ$, all pairs, and everything. The full list of classical and nonlinear options is in
Appendix~\ref{app:opponents}.}\label{tab:students}
\begin{tabular}{@{}p{2.6cm}p{7.6cm}p{3.2cm}@{}}\toprule
student & where its features come from & options to choose from\\\midrule
\textbf{Exact} & the exact quantum simulation & 9 groups $\times$ 5 couplings $\times$ 12 anneal times $=540$\\
\textbf{Mean-field} & the mean-field simulation & 540 (same menu)\\
\textbf{No coupling} & the $J=0$ simulation & $9\times12=108$\\
\textbf{Ground state} & the lowest-energy spin arrangement (the $\tau\to\infty$ end point) & 3 groups $\times$ 5 couplings $=15$\\
\textbf{Classical features} & generic feature expansions (Fourier, polynomials, splines, \ldots) and 22 tuned feature maps & $\sim$48\\
\textbf{Nonlinear} & the raw inputs, fed directly to gradient boosting, an RBF-kernel SVM, or a small neural net & 22 settings\\\bottomrule
\end{tabular}
\end{table}

\paragraph{Choosing and testing fairly (nested cross-validation).} The 600 samples are split into 5 outer
folds; each round, a student trains on 480 and is tested on the remaining 120. Inside the 480 it splits again
into 3 inner folds and uses them to pick its best option and its regularization strength $C$ (from $10^{-3}$ to
$10^{2}$). It then refits that choice on all 480 and is scored on the held-out 120. The student's score is the
average test AUC over the 5 rounds. Every student uses the \emph{same} folds, so comparisons are paired, and
methods that learn from the labels (such as the GAM shape fits) are refitted inside each fold, so the
held-out samples never leak into the choice.

\paragraph{The comparison.}
\begin{equation}
\text{margin}=\mathrm{AUC}_{\text{exact}}-\max\big(\mathrm{AUC}\ \text{of every other student}\big).
\end{equation}

\paragraph{The open-book student.} The exact student's menu always contains the exact recipe the teacher
used (same simulator, coupling, anneal time and measurements). On exact-teacher data it therefore acts as the
\emph{ceiling} (about 0.94 out of a possible 0.95), and what we are really measuring is whether anyone
\emph{else} can imitate it. This cuts both ways: on mean-field teacher data the mean-field student is the
open-book one, and the exact student cannot imitate \emph{it} either.

% =====================================================================

% =====================================================================
\section{What we found}

\textbf{Headline.} On labels the quantum map made itself, classical methods and the cheap non-quantum
versions of the map catch up almost everywhere. The exception is one region: strong coupling, a medium-to-long
anneal, and a readout that includes $\langle X_i\rangle$ and $\langle Y_i\rangle$. There, nothing we tried
catches up.

\begin{table}[H]\centering\small
\caption{Exact-teacher datasets by region. The \emph{margin} is how far the exact quantum student is ahead of
the \emph{best} opponent (in AUC), with all opponents included (classical feature models, tuned maps,
nonlinear learners, and the mean-field, no-coupling and ground-state versions of the map). $c$ is the coupling
level of Table~\ref{tab:J}.}\label{tab:regions}
\begin{tabular}{@{}p{3.3cm}p{3.6cm}cccc@{}}\toprule
region & settings & datasets & \shortstack{median\\margin} & \shortstack{margin\\$\ge0.05$} & \shortstack{catches\\up?}\\\midrule
$Z$-only readout & any $c$, any $\tau$ & 180 & 0.00 & 0\% & yes\\
weak coupling or short anneal & $c=0.5$, or $\tau\le1$ & 60 & 0.00 & 0\% & yes\\
in between & $c=1$ & 30 & 0.03 & 30\% & sometimes\\
edges & $c\ge2$ at $\tau=1.5$--2, or $\tau=40$ & 27 & 0.04 & 33\% & sometimes\\
\textbf{the window} & $\boldsymbol{c\ge2}$, $\boldsymbol{\tau=3}$--\textbf{20}, \textbf{full readout} & \textbf{63} & \textbf{0.15} & \textbf{87\%} & \textbf{no}\\\bottomrule
\end{tabular}
\end{table}

\textbf{1. Inside the window, nobody gets close.} The exact student scores about 0.94 (the ceiling is about
0.95). The best opponent reaches about 0.79, which is only about 65\% of the achievable signal above chance.
That holds against every opponent we tried: Fourier, polynomial and spline features (up to 128 per input),
22 tuned feature maps, gradient boosting, an RBF-kernel SVM and a small neural net on the raw inputs, and the
mean-field, no-coupling and ground-state versions of the map. Table~\ref{tab:window} shows where the window
sits in coupling and anneal time.

\begin{table}[H]\centering\small
\caption{Full readout: number of exact-teacher datasets (out of 3, one per $p$) where the exact student beats
every opponent by at least 0.05 AUC. Rows are coupling levels (Table~\ref{tab:J}), columns anneal times
$\tau$, with every opponent included (classical features, tuned maps, nonlinear learners, and the
mean-field, no-coupling and ground-state students). With $Z$-only readout every cell is 0.}\label{tab:window}
\begin{tabular}{@{}ccccccccccccc@{}}\toprule
$c$ $\backslash$ $\tau$ & 0.5 & 1 & 1.5 & 2 & 3 & 4 & 5 & 7 & 10 & 15 & 20 & 40\\\midrule
0.5 & 0 & 0 & 0 & 0 & 0 & 0 & 0 & 0 & 0 & 0 & 0 & 0\\
1   & 0 & 0 & 0 & 0 & 1 & 1 & 1 & 2 & 0 & 1 & 2 & 1\\
2   & 0 & 0 & 0 & 2 & \textbf{3} & \textbf{3} & \textbf{3} & \textbf{3} & \textbf{3} & \textbf{3} & \textbf{3} & 0\\
4   & 0 & 0 & 0 & 0 & 2 & \textbf{3} & \textbf{3} & \textbf{3} & \textbf{3} & \textbf{3} & \textbf{3} & 3\\
8   & 0 & 0 & 1 & 0 & 0 & 1 & 2 & 2 & \textbf{3} & \textbf{3} & \textbf{3} & 3\\\bottomrule
\end{tabular}
\end{table}

\textbf{2. Being hard for classical models isn't proof of anything quantum.} Labels made by the
\emph{mean-field} simulator, which is classical and cheap, are just as hard for the same opponents in that
region: they reach about 0.71, against 0.95 for the matching mean-field student. So ``it takes a lot of
Fourier or polynomial features to catch up'' is not, on its own, evidence for the quantum part. The fair test
is always the map's own classical version.

\textbf{3. Entanglement alone isn't enough.} Short anneals ($\tau=1$) at strong coupling are highly entangled
(about 0.98 on the Meyer--Wallach global entanglement, which runs from 0 for no entanglement to 1 for maximal; it is the average over qubits of $1-|\vec r_i|^2$, with $\vec r_i=(\langle X_i\rangle,\langle Y_i\rangle,\langle Z_i\rangle)$, equivalently the average single-qubit linear entropy \citep{Meyer2002,Brennen2003}), yet their labels are easy to copy (margin $\approx0$).
What matters is whether the label \emph{function} can be imitated, not whether the state is entangled.

\textbf{Still pending.} Whether the window survives larger systems ($p$ up to 12; running).

% =====================================================================
\section{Shot noise: does it survive realistic measurement?}

On hardware every measurement is an average over a finite number of repetitions (``shots''). We re-ran the
exact student with each measurement replaced by its estimate from $S$ shots,
\begin{equation}
\widehat{\langle P\rangle}=\frac{2k}{S}-1,\qquad k\sim\mathrm{Binomial}\Big(S,\tfrac{1+\langle P\rangle}{2}\Big),
\end{equation}
keeping the labels and all classical opponents noiseless (they are computed exactly on an ordinary computer).

\begin{table}[H]\centering\small
\caption{Shot noise, exact teachers in the window (63 datasets). Only the exact student gets noise.}\label{tab:shots}
\begin{tabular}{@{}lccc@{}}\toprule
shots per measurement & exact student AUC & median margin & margin $\ge0.05$\\\midrule
noiseless & 0.941 & 0.149 & 87\%\\
10{,}000 & 0.937 & 0.144 & 87\%\\
\textbf{1{,}000} & \textbf{0.925} & \textbf{0.124} & \textbf{83\%}\\
100 & 0.856 & 0.062 & 56\%\\\bottomrule
\end{tabular}
\end{table}

The window is essentially intact at 1{,}000--10{,}000 shots and halves at 100, where the exact student itself
becomes too noisy. Even at 100 shots the core of the window (coupling levels 2--4, $\tau=4$--10) still wins on
most datasets. Outside the window nothing changes. The $\tau=40$ cells, whose advantage rested on very small
$\langle X_i\rangle/\langle Y_i\rangle$ signals, lose it entirely at 1{,}000 shots, as expected
(Table~\ref{tab:window1000}).

\begin{table}[H]\centering\small
\caption{As Table~\ref{tab:window}, but with the exact student's measurements estimated from 1{,}000 shots.}\label{tab:window1000}
\begin{tabular}{@{}ccccccccccccc@{}}\toprule
$c$ $\backslash$ $\tau$ & 0.5 & 1 & 1.5 & 2 & 3 & 4 & 5 & 7 & 10 & 15 & 20 & 40\\\midrule
0.5 & 0 & 0 & 0 & 0 & 0 & 0 & 0 & 0 & 0 & 0 & 0 & 0\\
1   & 0 & 0 & 0 & 0 & 1 & 1 & 1 & 2 & 0 & 1 & 1 & 0\\
2   & 0 & 0 & 0 & 0 & \textbf{3} & \textbf{3} & \textbf{3} & \textbf{3} & \textbf{3} & \textbf{3} & \textbf{3} & 0\\
4   & 0 & 0 & 0 & 0 & 2 & \textbf{3} & \textbf{3} & \textbf{3} & \textbf{3} & 2 & \textbf{3} & 0\\
8   & 0 & 0 & 0 & 0 & 0 & 0 & 2 & 2 & \textbf{3} & \textbf{3} & 2 & 0\\\bottomrule
\end{tabular}
\end{table}

A practical note: ``1{,}000 shots'' is per measurement basis per sample. The full readout needs three bases
($X$, $Y$, $Z$), so about 3{,}000 circuit runs per sample, or about $1.8$ million per (coupling, $\tau$)
setting for 600 samples.

% =====================================================================
% Planned sections (to be written together, one at a time):


% \section{Scope}
% \section{Next steps}            spotting window-like datasets beforehand; scaling run
% =====================================================================

\bibliographystyle{plainnat}
\begin{thebibliography}{10}\small
% TODO: verify first-author initial and author list
\bibitem[Simen et~al.(2025)]{Simen2025}
Simen et~al. Digitized counterdiabatic quantum feature extraction. arXiv:2510.13807, 2025.

\bibitem[Liu et~al.(2021)]{Liu2021}
Y.~Liu, S.~Arunachalam, K.~Temme. A rigorous and robust quantum speed-up in supervised machine learning.
\emph{Nature Physics}, 2021. \href{https://arxiv.org/abs/2010.02174}{arXiv:2010.02174}.

\bibitem[Fujii and Nakajima(2017)]{Fujii2017}
K.~Fujii, K.~Nakajima. Harnessing disordered-ensemble quantum dynamics for machine learning.
\emph{Phys.\ Rev.\ Applied}, 2017. \href{https://arxiv.org/abs/1602.08159}{arXiv:1602.08159}.

% TODO: verify
\bibitem[Schuld et~al.(2021)]{Schuld2021}
M.~Schuld, R.~Sweke, J.~J.~Meyer. Effect of data encoding on the expressive power of variational quantum
machine learning models. \emph{Phys.\ Rev.\ A} 103, 032430, 2021.

\bibitem[Schreiber et~al.(2023)]{Schreiber2023}
F.~J.~Schreiber, J.~Eisert, J.~J.~Meyer. Classical surrogates for quantum learning models.
\emph{Phys.\ Rev.\ Lett.} 131, 100803, 2023. \href{https://link.aps.org/doi/10.1103/PhysRevLett.131.100803}{doi}.

\bibitem[Landman et~al.(2022)]{Landman2022}
J.~Landman, S.~Thabet, C.~Dalyac, H.~Mhiri, E.~Kashefi. Classically approximating variational quantum machine
learning with random Fourier features. \href{https://arxiv.org/abs/2210.13200}{arXiv:2210.13200}, 2022.

\bibitem[Sweke et~al.(2025)]{Sweke2025}
R.~Sweke et~al. Potential and limitations of random Fourier features for dequantizing quantum machine
learning. \emph{Quantum}, 2025. \href{https://quantum-journal.org/papers/q-2025-02-20-1640/}{link}.

\bibitem[Huang et~al.(2021)]{Huang2021}
H.-Y.~Huang et~al. Power of data in quantum machine learning. \emph{Nature Communications}, 2021.
\href{https://arxiv.org/abs/2011.01938}{arXiv:2011.01938}.

\bibitem[Bowles et~al.(2024)]{Bowles2024}
J.~Bowles, S.~Ahmed, M.~Schuld. Better than classical? The subtle art of benchmarking quantum machine learning
models. \href{https://arxiv.org/abs/2403.07059}{arXiv:2403.07059}, 2024.

\bibitem[Meyer and Wallach(2002)]{Meyer2002}
D.~A.~Meyer, N.~R.~Wallach. Global entanglement in multiparticle systems. \emph{J.\ Math.\ Phys.}, 2002.
\href{https://arxiv.org/abs/quant-ph/0108104}{arXiv:quant-ph/0108104}.

\bibitem[Brennen(2003)]{Brennen2003}
G.~K.~Brennen. An observable measure of entanglement for pure states of multi-qubit systems.
\emph{Quantum Inf.\ Comput.}, 2003. \href{https://arxiv.org/abs/quant-ph/0305094}{arXiv:quant-ph/0305094}.

\bibitem[Cerezo et~al.(2023)]{Cerezo2023}
M.~Cerezo et~al. Does provable absence of barren plateaus imply classical simulability?
\href{https://arxiv.org/abs/2312.09121}{arXiv:2312.09121}, 2023.
\end{thebibliography}

% =====================================================================
\appendix
\section{Full list of opponents}\label{app:opponents}

Throughout, $q=p(p-1)/2$ is the number of qubit pairs, inputs are the standardized $z$, and every option is
followed by logistic regression with $C$ chosen by inner cross-validation unless stated otherwise.

\paragraph{Generic feature expansions (no labels used to build them).} Each size is also scored on its own
to give the ``how many features does it take'' curves.
\begin{itemize}\setlength{\itemsep}{1pt}
\item \textbf{Fourier:} $\sin(kz_j)$, $\cos(kz_j)$ for increasing $k$, filled evenly across features, with
$D=m\,p$ features for $m\in\{1,2,4,8,16,32,64,128\}$.
\item \textbf{Random Fourier features} (approximating an RBF kernel, $\gamma=1/p$) at the same sizes.
\item \textbf{Polynomials} of degree 1--6 (all products up to that degree), kept while $D\le128p$.
\item \textbf{Cubic splines} per feature with 3, 5, 9, 17, 33 or 65 knots.
\end{itemize}

\paragraph{The 22 tuned feature maps.}
\begin{itemize}\setlength{\itemsep}{1pt}
\item \textbf{Raw} inputs ($p$ features).
\item \textbf{Balanced Fourier} and \textbf{random Fourier} at three fixed sizes, $p$, $3p$ and $3p+3q$
(6 maps).
\item \textbf{Hermite polynomials:} per feature (degrees 1--3), pairwise products (degree 1, and degrees
1--3), and both together (4 maps).
\item \textbf{Quadratic splines} per feature (2 knots) and \textbf{pairwise interactions} $z_iz_j$ with the raw
inputs (2 maps).
\item \textbf{Label-matched Fourier:} for each feature, the frequencies most correlated with the label
(periodogram peaks), with $p$ or $3p$ features (2 maps; learned from the training labels).
\item \textbf{GAM shapes:} a smooth function of each feature, fitted to the labels; their pairwise products;
and both together (3 maps; learned from the labels).
\item \textbf{GA2M:} per-feature shapes plus pairwise interaction shapes fitted jointly to the labels; the pair
part alone, and everything (2 maps; learned from the labels).
\item \textbf{Symmetric aggregates:} splines of $\mathrm{mean}(z)$ and $\mathrm{mean}(z^2)$ across features,
alone and appended to the raw inputs (2 maps).
\end{itemize}
Maps marked ``learned from the labels'' are refitted inside every cross-validation fold. The
\emph{classical features} student picks the best of all generic sizes and the 22 maps by inner
cross-validation.

\paragraph{Nonlinear learners on the raw inputs} (22 settings, best per learner and overall chosen by inner
cross-validation):
\begin{itemize}\setlength{\itemsep}{1pt}
\item \textbf{Gradient boosting} (histogram-based, 300 rounds): learning rate $\in\{0.05,0.1\}$, maximum
leaves per tree $\in\{8,31\}$ (4 settings).
\item \textbf{RBF-kernel SVM:} $C\in\{0.3,1,3,10,30\}$, $\gamma\in\{0.3,1,3\}/p$ (15 settings).
\item \textbf{Neural net:} two hidden layers of 64 units, weight penalty $\in\{10^{-4},10^{-2},1\}$, early
stopping (3 settings).
\end{itemize}

\paragraph{The map's own students.} Exact, mean-field and no-coupling students choose among the 9 measurement
groups ($X$, $Y$, $Z$, $XYZ$, $XX$, $YY$, $ZZ$, all pairs, everything), the 5 couplings of
Table~\ref{tab:J} (not for $J=0$) and the 12 anneal times. The ground-state student chooses among the spins
$\sigma_i$, the pair products $\sigma_i\sigma_j$, or both, at each of the 5 couplings.



\end{document}
